% Einheit 1 von 3 – Satz und Hypotenuse (bank/pythagoras/e1.jsonl, zusammenbau v0.1)
%% TODO Zweigzeile Teil 1: was hier gelernt wird (Satz in Schülersprache aus dem Einheitstitel) – Einheit 1
\einheitenkopf[e1]{Einheit 1 von 3 · Satz und Hypotenuse}
\zweigzeile{ab Kl. 8, je nach Buch bis Kl. 9 · P10 oft}
\setcounter{aufgabe}{11}

%% TODO Titel: Ich-kann-Satz fehlt in der Bank (Platzhalter: Kettenname) – Nr. 12
%% TODO Anweisung: ein Satz über der ersten Teilaufgabe fehlt in der Bank – Nr. 12
\begin{aufgabe}{Gilt der Satz hier?}
\begin{teile}
\teil Gilt der Satz des Pythagoras? \janein
\end{teile}
\end{aufgabe}

%% TODO Titel: Ich-kann-Satz fehlt in der Bank (Platzhalter: Kettenname) – Nr. 13
%% TODO Anweisung: ein Satz über der ersten Teilaufgabe fehlt in der Bank – Nr. 13
\begin{aufgabe}{Hypotenuse}
\swa{Rechten Winkel einkreisen, Hypotenuse mit H – wie lang ist sie? \leerfeld[cm]}{\dreieckrw{5}{1.12}{$40$ cm}{$9$ cm}{$41$ cm}}
%% TODO Darstellung für schwach fehlt in der Bank (pythagoras-e1-k2-s1-v1; Abschnitt „Für schwache Schüler“)
\swz{Katheten $a = 30$ cm, $b = 40$ cm – Hypotenuse $c$?}{}
%% TODO Darstellung für schwach fehlt in der Bank (pythagoras-e1-k2-s1-v2; Abschnitt „Für schwache Schüler“)
\swz{Katheten $a = 9$ m, $b = 40$ m – Hypotenuse $c$?}{}
%% TODO Darstellung für schwach fehlt in der Bank (pythagoras-e1-k2-s1-v3; Abschnitt „Für schwache Schüler“)
\swz{Katheten $a = 18$ cm, $b = 24$ cm – Hypotenuse $c$?}{}
%% TODO Darstellung für schwach fehlt in der Bank (pythagoras-e1-k2-s1-v4; Abschnitt „Für schwache Schüler“)
\swz{Katheten $a = 11$ mm, $b = 60$ mm – Hypotenuse $c$?}{}
%% TODO Darstellung für schwach fehlt in der Bank (pythagoras-e1-k2-s1-v5; Abschnitt „Für schwache Schüler“)
\swz{Katheten $a = 40$ m, $b = 42$ m – Hypotenuse $c$?}{}
\swfrage{Was bleibt gleich, was ändert sich?}
\swz{Wie lang ist die mit ? markierte Seite?}{\dreieck{(0,0)}{(5,0)}{(1.4,2.24)}{$45$ cm}{$28$ cm}{?}{}{}{$90^\circ$}}
%% TODO Darstellung für schwach fehlt in der Bank (pythagoras-e1-k2-s3-v1; Abschnitt „Für schwache Schüler“)
\swz{Katheten $a = 4$ cm, $b = 9$ cm – Hypotenuse $c$ auf eine Dezimale?}{}
%% TODO Darstellung für schwach fehlt in der Bank (pythagoras-e1-k2-s4-v1; Abschnitt „Für schwache Schüler“)
\swz{Katheten $a = 2{,}5$ m, $b = 6$ m – Hypotenuse $c$?}{}
%% TODO Darstellung für schwach fehlt in der Bank (pythagoras-e1-k2-s5-v1; Abschnitt „Für schwache Schüler“)
\swz{Katheten $a = 1{,}2$ m, $b = 50$ cm – Hypotenuse $c$ in cm?}{}
\swz{Gleichung nach dem Satz des Pythagoras?}{\dreieckrw{5}{3.57}{$m$}{$n$}{$k$}}
\begin{teile}
\teil Welche Aussage gilt im rechtwinkligen Dreieck? \\ \kreuz{Die Hypotenuse ist die kürzeste Seite.} \\ \kreuz{Das Quadrat über der Hypotenuse ist so groß wie die zwei Quadrate über den Katheten zusammen.} \\ \kreuz{Die zwei Katheten sind zusammen so lang wie die Hypotenuse.}
\end{teile}
\begin{teile}
\teil Welche Gleichung gilt für $g$? \\ \kreuz{$g = \sqrt{e^2 + f^2}$} \\ \kreuz{$g = e^2 + f^2$} \\ \kreuz{$g = \sqrt{e^2 - f^2}$} \\ \kreuz{$g = e + f$}
\end{teile}
%% TODO Darstellung für schwach fehlt in der Bank (pythagoras-e1-k2-s9-v1; Abschnitt „Für schwache Schüler“)
\swz{Rechteck $45$ cm mal $24$ cm – Diagonale?}{}
\swz{Eine rechteckige Wiese ist $60$ m lang und $11$ m breit. Wie lang ist der Weg quer über die Wiese von Ecke zu Ecke?}{\dreieckrw{5}{0.92}{$60$ m}{$11$ m}{?}}
\swz{Eine Leiter steht $0{,}9$ m von einer Mauer entfernt und reicht genau bis zur Mauerkante in $4$ m Höhe. Oben ragt sie noch $1$ m über die Kante hinaus. Wie lang ist die Leiter?}{\dreieckrw{1.12}{5}{$0{,}9$ m}{$4$ m}{?}}
\swz[3]{Ein dreieckiges Grundstück ABC hat bei C einen rechten Winkel. Der Zaun von B nach C ist $14$ m lang, der von A nach C $39$ m. Berechne die Länge der Seite AB. Gib das Zwischenergebnis $AB^2$ an und runde auf eine Dezimale \hfill (P10 2020 OS).}{\dreieckrw{5}{1.79}{$39$ m}{$14$ m}{?}}
\swz[3]{Eine Laderampe führt auf einer waagerechten Länge von $240$ cm um $18$ cm nach oben. Berechne die Länge der Rampe auf eine Dezimale \hfill (P10 2025 OS).}{\dreieckrw{6}{0.45}{$240$ cm}{$18$ cm}{?}}
\end{aufgabe}

%% TODO Titel: Ich-kann-Satz fehlt in der Bank (Platzhalter: Kettenname) – Nr. 14
%% TODO Anweisung: ein Satz über der ersten Teilaufgabe fehlt in der Bank – Nr. 14
\begin{aufgabe}{Quadrate über den Seiten zeichnen und Flächen vergleichen (Kästchen zählen)}
\swa{Dreieck ABC, rechter Winkel bei A. Zeichne die Quadrate über allen drei Seiten. Wie viele Kästchen hat jedes? \leerfeld + \leerfeld = \leerfeld}{\begin{ksys}[xmin=0,xmax=10,ymin=0,ymax=9] \punkt{4}{3}{A} \punkt{6}{3}{B} \punkt{4}{6}{C} \end{ksys}}
\end{aufgabe}

%% TODO Titel: Ich-kann-Satz fehlt in der Bank (Platzhalter: Kettenname) – Nr. 15
%% TODO Anweisung: ein Satz über der ersten Teilaufgabe fehlt in der Bank – Nr. 15
\begin{aufgabe}{Ergebnis mit sinnvoller Genauigkeit angeben}
%% TODO Darstellung für schwach fehlt in der Bank (pythagoras-e1-k4-s1-v1; Abschnitt „Für schwache Schüler“)
\swz{Katheten $2{,}3$ m und $3{,}1$ m – Hypotenuse $c$, sinnvoll gerundet?}{}
\end{aufgabe}

%% TODO Titel: Ich-kann-Satz fehlt in der Bank (Platzhalter: Kettenname) – Nr. 16
%% TODO Anweisung: ein Satz über der ersten Teilaufgabe fehlt in der Bank – Nr. 16
\begin{aufgabe}{Hypotenuse – Fehler finden, Begründen, Anwendung, Darstellungswechsel}
\swz[3]{Katheten $33$ cm und $56$ cm. Lara rechnet: \rechnung{c^2 &= 33^2 + 56^2 \\ c^2 &= 4\,225 \\ c &= 4\,225 \text{ cm}} Wo ist der Fehler?}{}
\swfrage{Warum ist die Hypotenuse länger als jede Kathete?}
\swz[3]{Ein Fernseher hat einen Bildschirm, der $111$ cm breit und $62$ cm hoch ist. Im Laden steht er als „50 Zoll“ ($1$ Zoll $= 2{,}54$ cm, gemessen wird die Diagonale). Stimmt die Angabe?}{}
\swa{Eine Leiter lehnt an einer Wand; ihr Fuß steht $1{,}5$ m von der Wand entfernt. Zeichne eine Skizze, markiere den rechten Winkel und beschrifte die Hypotenuse mit H.}{\rechenplatz[halb]{4}}
\end{aufgabe}

\uebersichtskasten{Hypotenuse: Die Seite gegenüber dem rechten Winkel heißt Hypotenuse und ist die längste Seite; die beiden Seiten am rechten Winkel heißen Katheten. \\ Satz des Pythagoras: Im rechtwinkligen Dreieck ist das Quadrat über der Hypotenuse so groß wie die beiden Quadrate über den Katheten zusammen – Katheten hoch zwei addieren, dann die Wurzel ziehen: a² + b² = c², c = √(a² + b²). \\ a = 6 cm, b = 8 cm: c² = 6² + 8² = 36 + 64 = 100 \quad c = √100 = 10 cm \\ Geht die Wurzel nicht auf, bleibt der gerundete Taschenrechnerwert stehen. \quad a = 5 cm, b = 7 cm: c² = 25 + 49 = 74 \quad c = √74 ≈ 8,6 cm \\ Der Buchstabe der Hypotenuse steht allein auf einer Seite der Gleichung, egal wie er heißt: Katheten x und y, Hypotenuse z → x² + y² = z².}
